% !PRINTED 467
\weberpage{467}

Now $G(t)$ and $g(t)$ have a root in common. But since $g(t)$ is irreducible, $G(t)$ must be divisible by $g(t)$; hence all the roots of $g(t)$ are also contained in $N$, i.\,e. $N$ is a normal field, q.\,e.\,d.

Every equation $g(t) = 0$ is called a \gesperrt{Galois resolvent} of the equation $F(x) = 0$, and a \gesperrt{Galois resolvent} is accordingly defined by the following specification:

\gesperrt{An equation $g(t) = 0$ is a Galois resolvent of a given equation $F(x) = 0$ in $\Omega$ if 1) $g(t)$ is irreducible, if 2) all the roots of $F(x)$ are rationally expressible through a root $\varrho$ of $g(t)$, and 3) one root of $g(t)$ is rationally expressible through the roots of $F(x)$.}

For by 2), $N$ is contained in $\Omega(\varrho)$, and by 3) one of the fields conjugate to $\Omega(\varrho)$, namely $\Omega(\varrho_1)$, is contained in $N$. The degrees of $N$ and $\Omega(\varrho)$ therefore cannot be different, and consequently $N = \Omega(\varrho)$.

\gesperrt{Every Galois resolvent is a normal equation, and a normal equation is its own Galois resolvent.}

\begin{center}
\S.\,146.

\gesperrt{The Substitutions of a Normal Field.}
\end{center}

Now let $\Omega(\varrho)$ be a normal field of degree $\mu$ and $\varrho$ one of its primitive numbers, and further let
\weq{1}{g(t) = 0}
be the irreducible equation of degree $\mu$ of which one root is $t = \varrho$. The elements conjugate to $\varrho$ are as follows.
\weq{2}{\varrho,\ \varrho_1,\ \varrho_2 \ldots \varrho_{\mu-1}.}
Since by the definition of the normal field the quantities (2) are all contained in $\Omega(\varrho)$, we can determine $\Theta_1(t), \Theta_2(t) \ldots \Theta_{\mu-1}(t)$ as integral rational functions in $\Omega$, of degree at most $\mu - 1$, such that
\weq{3}{\varrho_1 = \Theta_1(\varrho),\ \varrho_2 = \Theta_2(\varrho) \ldots \varrho_{\mu-1} = \Theta_{\mu-1}(\varrho)}
%CONT>>
